\chapter{Normalizzazione Relazionale}

Nel capitolo precedente abbiamo affrontato la transizione dal modello concettuale (Class Diagram UML) agli schemi relazionali. Durante la fase di ristrutturazione (si vedano, ad esempio, i paragrafi sulla risoluzione degli attributi strutturati e multivalore), abbiamo già sfiorato e applicato intuitivamente alcuni principi volti a "pulire" le tabelle e separare i concetti. 

Tuttavia, l'ingegneria del software non può basarsi unicamente sulle euristiche o sull'intuito del progettista. In questo capitolo tratteremo la normalizzazione (e le relative Forme Normali) in modo rigoroso, dettagliato e formale. Questa teoria rappresenta il "collaudo" definitivo degli schemi relazionali appena prodotti: fornisce gli strumenti esatti per dimostrare che la struttura del database è solida e immune da difetti. Solo superando questo test formale, il discorso sulla progettazione logica può considerarsi completato e pronto per la traduzione in codice SQL.

\section{Ridondanze e Anomalie: Il Problema Strutturale}

Un database relazionale è progettato male quando un singolo schema (una tabella) raggruppa forzatamente concetti eterogenei che non dovrebbero stare insieme. Questo errore di modellazione genera \textbf{ridondanza} (la duplicazione inutile della stessa informazione su più righe) che, a sua volta, innesca tre tipologie di anomalie critiche durante il normale ciclo di vita del software.

\begin{example}[Il disastro della tabella non normalizzata]
Supponiamo di aver modellato un'unica tabella mista per gestire i dipendenti e i dipartimenti in cui lavorano, eleggendo la matricola a chiave primaria:

\texttt{DIPENDENTE(\underline{matricola}, nome, dipartimento, sede\_dipartimento)}

Se il dipartimento "IT" si trova a "Milano" e ha 1000 dipendenti, la stringa "Milano" verrà memorizzata fisicamente 1000 volte. Questa ridondanza apre le porte alle tre anomalie relazionali:

\begin{enumerate}
    \item \textbf{Anomalia di Aggiornamento (Update Anomaly):} 
    Se il dipartimento "IT" viene trasferito a "Roma", l'applicativo dovrà eseguire una \texttt{UPDATE} su 1000 righe distinte. Se per un calo di rete o un crash del server l'aggiornamento si interrompe a 500 righe, il database entra in uno stato \textit{inconsistente}: metà dei dipendenti IT risulterà a Roma, l'altra metà a Milano. Il dato perde affidabilità.
    
    \item \textbf{Anomalia di Inserimento (Insert Anomaly):}
    L'azienda apre un nuovo dipartimento "Marketing" a "Napoli", ma non ha ancora assunto nessun dipendente. Poiché la \texttt{matricola} è chiave primaria (quindi \texttt{NOT NULL} per definizione), il database rifiuterà l'inserimento. È strutturalmente impossibile registrare l'esistenza di un dipartimento vuoto.
    
    \item \textbf{Anomalia di Cancellazione (Delete Anomaly):}
    Il dipartimento "HR" ha un solo dipendente. Se questo dipendente si dimette, l'applicativo eseguirà una \texttt{DELETE} sulla sua riga. Oltre a cancellare il dipendente, \textbf{distruggeremo involontariamente anche l'informazione che il dipartimento HR esiste} e che si trovava in una determinata sede. Abbiamo perso un dato vitale per l'azienda a causa dell'eliminazione di un dato non correlato.
\end{enumerate}
\end{example}

Per risolvere alla radice queste anomalie, il progettista non deve scrivere codice applicativo più complesso, ma deve modificare la struttura del database. La soluzione consiste nello \textbf{spaccare} (decomporre) la tabella difettosa in due tabelle distinte, legate da una chiave esterna:
\begin{itemize}
    \item \texttt{DIPENDENTE(\underline{matricola}, nome, dipartimento*)}
    \item \texttt{DIPARTIMENTO(\underline{nome\_dip}, sede\_dipartimento)}
\end{itemize}
La teoria della normalizzazione (le Forme Normali) fornisce le regole esatte per capire \textit{quando} e \textit{come} effettuare questa scomposizione in modo sicuro.


\section{Dipendenze Funzionali}

Nel paragrafo precedente abbiamo visto come le ridondanze generino bug strutturali (le anomalie). Per evitare di dover scovare queste ridondanze "a occhio" o basandosi solo sull'intuito, la teoria relazionale fornisce uno strumento analitico per identificare le anomalie: la \textbf{Dipendenza Funzionale} (Spesso abbreviata in \textit{FD - Functional Dependency}).

Una dipendenza funzionale stabilisce che il valore di un determinato campo "decide" in modo univoco il valore di un altro campo. Se conosco il primo, non ho dubbi sul secondo.

Formalmente, data una relazione (tabella) $r$ definita su uno schema $R(X)$, e dati due sottoinsiemi di attributi non vuoti $Y$ e $Z$ (appartenenti a $X$), diremo che esiste su $r$ una dipendenza funzionale tra $Y$ e $Z$ se:

\begin{quote}
Per ogni coppia di tuple (righe) $t_1$ e $t_2$ presenti in $r$ che hanno \textbf{gli stessi valori sugli attributi $Y$}, risulta obbligatoriamente che $t_1$ e $t_2$ hanno \textbf{gli stessi valori anche sugli attributi $Z$}.
\end{quote}

In notazione algebrica, questo vincolo di integrità si scrive:
$$ Y \rightarrow Z $$
E si legge: \textit{"Y determina funzionalmente Z"}, oppure \textit{"Z dipende funzionalmente da Y"}. L'insieme $Y$ prende il nome di \textbf{Determinante}.

\begin{example}
Immaginiamo una tabella temporanea che traccia i veicoli assicurati e le relative polizze:

\begin{center}
    \texttt{COPERTURA(numero\_polizza, targa, marca, modello)}
\end{center}


Osservando la natura dei dati (ovvero la semantica del mondo reale), possiamo affermare con certezza matematica che esiste la dipendenza funzionale:
$$ \texttt{targa} \rightarrow \texttt{marca, modello} $$
Il motivo risponde esattamente alla definizione formale: se prendiamo due righe qualsiasi della tabella ($t_1$ e $t_2$) che presentano lo stesso valore nel campo \texttt{targa} (es. "AB123CD"), quelle due righe dovranno \textbf{obbligatoriamente} avere lo stesso valore nei campi \texttt{marca} (es. "Fiat") e \texttt{modello} (es. "Panda"). Un veicolo non può cambiare marca da una riga all'altra.
\end{example}

\paragraph{}

L'analisi delle dipendenze funzionali si basa su tre postulati pratici che semplificano il lavoro di progettazione:

\begin{enumerate}
    \item \textbf{Scomposizione (Decomposizione del conseguente):}
    Se un determinante $Y$ definisce un insieme di attributi $Z$ composto da $A_1, A_2, \dots, A_k$, allora possiamo affermare che la relazione soddisfa $Y \rightarrow Z$ se e solo se soddisfa singolarmente tutte le dipendenze: 
    $$ Y \rightarrow A_1 \quad \text{e} \quad Y \rightarrow A_2 \quad \text{e} \dots \quad Y \rightarrow A_k $$
    \textit{Esempio:} Dire che $\texttt{targa} \rightarrow \texttt{marca, modello}$ equivale a dire che $\texttt{targa} \rightarrow \texttt{marca}$ e $\texttt{targa} \rightarrow \texttt{modello}$. Per semplicità, d'ora in avanti considereremo sempre dipendenze del tipo $Y \rightarrow A$ (dove la parte destra è un singolo attributo).

    \item \textbf{Dipendenze Banali (da ignorare):}
    Se l'attributo $A$ fa già parte dell'insieme $Y$, la dipendenza viene definita \textbf{banale}. Ad esempio, la dipendenza $\texttt{targa, modello} \rightarrow \texttt{targa}$ è ovvia e non aggiunge alcuna informazione utile al design del database. Nello studio della normalizzazione, si prendono in esame esclusivamente le dipendenze funzionali \textbf{non banali}.

    \item \textbf{La generalizzazione del concetto di Chiave:}
    Cos'è una Chiave Primaria in ottica matematica? Non è altro che un determinante supremo. Dire che \texttt{numero\_polizza} è la chiave della tabella, equivale a dichiarare la dipendenza funzionale:
    $$ \texttt{numero\_polizza} \rightarrow \text{\textit{Tutti gli altri attributi dello schema}} $$
    La Dipendenza Funzionale, quindi, \textbf{generalizza il concetto di chiave}. 
\end{enumerate}

\begin{tip}[L'intuizione dietro le Forme Normali]
Questa generalizzazione ci fornisce la chiave di lettura per l'intera teoria della normalizzazione. Se in una tabella troviamo una dipendenza funzionale forte (come $\texttt{targa} \rightarrow \texttt{marca}$), ma la \texttt{targa} \textbf{non è} la chiave primaria di quella tabella, significa che c'è un difetto strutturale: stiamo costringendo un sotto-tema (il veicolo) dentro un tema più grande (la copertura). È esattamente questa discrepanza a generare le ridondanze e le anomalie che impongono la scomposizione in più tabelle.
\end{tip}

\section{La Forma Normale di Boyce-Codd (BCNF)}

Il traguardo ultimo della progettazione logica relazionale è garantire che il database sia totalmente privo di anomalie e ridondanze. Lo strumento che certifica questa garanzia strutturale è la Forma Normale di Boyce-Codd (BCNF). 

\paragraph{Definizione della BCNF}
Uno schema relazionale $R$ è in Forma Normale di Boyce-Codd se, per ogni sua dipendenza funzionale non banale $X \rightarrow Y$, risulta che \textbf{$X$ è una superchiave} per lo schema $R$.

In termini pratici e ingegneristici, questa regola si traduce in un diktat assoluto: \textit{"L'unica cosa che ha il diritto di determinare il valore di un attributo è la chiave primaria (o una chiave candidata) della tabella stessa."} Se un attributo non-chiave si arroga il diritto di determinare altri attributi, il database nasconde un difetto strutturale.

\begin{example}[Violazione della BCNF e Ritorno alle Anomalie]
Riprendiamo lo schema difettoso analizzato nel primo paragrafo e applichiamo la lente delle dipendenze funzionali per dimostrarne l'invalidità:

$$ \texttt{DIPENDENTE(\underline{matricola}, nome, dipartimento, sede\_dipartimento)} $$

In questo schema, l'unica chiave candidata (eletta a chiave primaria) è \texttt{matricola}. Di conseguenza, la \texttt{matricola} è una superchiave. La dipendenza funzionale $\texttt{matricola} \rightarrow \texttt{nome, dipartimento, sede\_dipartimento}$ è quindi del tutto lecita e rispetta la regola della BCNF.

Tuttavia, analizzando la semantica del dominio aziendale, sappiamo che un dipartimento si trova in una specifica sede. Esiste quindi una seconda dipendenza funzionale non banale "nascosta" nello schema:
$$ \texttt{dipartimento} \rightarrow \texttt{sede\_dipartimento} $$

Sottoponiamo questa dipendenza al test della BCNF: il determinante \texttt{dipartimento} è una superchiave dello schema \texttt{DIPENDENTE}? \textbf{Assolutamente no.} Il nome di un dipartimento non identifica in modo univoco una riga, poiché possono coesistere centinaia di impiegati assegnati allo stesso dipartimento.

Poiché esiste almeno una dipendenza funzionale il cui determinante non è una superchiave, l'intero schema \textbf{viola la BCNF}. È esattamente questa specifica violazione a generare le ridondanze e le letali anomalie di inserimento, aggiornamento e cancellazione.
\end{example}

\paragraph{La Decomposizione: Il ripristino della BCNF}
Quando uno schema fallisce il collaudo della BCNF, la teoria relazionale impone di decomporlo. Il processo consiste nell'isolare la dipendenza funzionale "illegale" e promuoverla a tabella autonoma. In questo nuovo schema, il vecchio determinante difettoso diventa finalmente una chiave primaria legittima.

Applicando la scomposizione al nostro esempio, otteniamo due nuovi schemi:
\begin{itemize}
    \item \texttt{DIPARTIMENTO(\underline{nome\_dip}, sede\_dipartimento)}
    \begin{itemize}
        \item Dipendenze presenti: $\texttt{nome\_dip} \rightarrow \texttt{sede\_dipartimento}$.
        \item Test BCNF: \texttt{nome\_dip} è la chiave primaria. L'unica dipendenza rispetta la BCNF.
    \end{itemize}
    \item \texttt{DIPENDENTE(\underline{matricola}, nome, dipartimento*)}
    \begin{itemize}
        \item Dipendenze presenti: $\texttt{matricola} \rightarrow \texttt{nome, dipartimento}$.
        \item Test BCNF: \texttt{matricola} è la chiave primaria. L'unica dipendenza rispetta la BCNF.
    \end{itemize}
\end{itemize}

Entrambi gli schemi generati si trovano ora rigorosamente in Forma Normale di Boyce-Codd. Il database è strutturalmente inattaccabile.